\documentclass[10pt,twocolumn,letterpaper]{article}
\usepackage[review]{cvpr}
\usepackage{times}
\usepackage{graphicx}
\usepackage{amsmath}
\usepackage{amssymb}
\usepackage{booktabs}
\usepackage{multirow}
\usepackage{xcolor}
\usepackage[pagebackref,breaklinks,colorlinks]{hyperref}

\def\cvprPaperID{****}
\def\confName{CVPR}
\def\confYear{2025}

\begin{document}

\title{Depth Super-Resolution for 3D Gaussian Splatting Face Reconstruction:\\Rendering-Induced Degradation and Zone-Aware Evaluation}

\author{Anonymous}

\maketitle

% ===================================================================
\begin{abstract}
We address the task of depth map super-resolution in the context of 3D Gaussian Splatting (3DGS) face reconstruction.
Given a 256$\times$256, 8-bit depth map rendered by a pre-trained 3DGS model, our goal is to recover a 1024$\times$1024, 16-bit high-fidelity depth map---jointly performing $\times$4 spatial upsampling and 8-bit to 16-bit dequantization.

We make three contributions.
First, we identify \emph{rendering-induced degradation} as a distinct depth super-resolution setting: the combination of splat-boundary artifacts, opacity fall-off, and 8-bit quantization produced by 3DGS rendering differs fundamentally from the bicubic downsampling assumed by prior depth SR methods.
A 5$\times$5 cross-degradation test shows a 3.5\,dB PSNR gap between matched and mismatched train/test degradations, confirming that models must be trained on the target degradation type.
Second, we propose \emph{zone-aware evaluation}, restricting depth SR metrics to the near-frontal confidence cone ($|\text{yaw}| \leq 20^\circ$) where 3DGS geometry is reliable, avoiding contamination from side/back-view hallucinations.
Third, we show that 3DGS-rendered surface normals are unusable as a supervision signal due to per-splat aliasing, while DSINE pseudo-ground-truth normals~\cite{bae2024dsine} provide a clean alternative that matches baseline performance (L$_1$ = 0.00228 \vs 0.00232).

Experiments on 1{,}287 FFHQ-derived 3DGS face scans show that a simple UNet with bicubic residual learning achieves 46.2\,dB PSNR, outperforming CVPR~2025 SOTA DORNet~\cite{yan2025dornet} by 4.1\,dB and AAAI~2024 SGNet~\cite{yan2024sgnet} by 3.1\,dB when these are applied zero-shot or trained on the same data.
\end{abstract}

% ===================================================================
\section{Introduction}
\label{sec:intro}

3D Gaussian Splatting (3DGS)~\cite{kerbl20233dgs} has emerged as a powerful representation for real-time novel view synthesis, with recent extensions to single-image 3D face reconstruction such as FaceLift~\cite{wu2024facelift}.
While FaceLift produces plausible multi-view RGB and depth maps from a single photograph, the rendered depth maps suffer from two compounding limitations: (1)~limited spatial resolution, typically 256$\times$256 constrained by the input image, and (2)~8-bit quantization that discards fine geometric detail.

Depth super-resolution (SR) is a natural remedy, yet applying existing depth SR methods trained on standard benchmarks (\eg, NYU-v2~\cite{silberman2012nyuv2}, Middlebury~\cite{scharstein2014middlebury}) yields poor results on 3DGS-rendered depth.
The reason is a fundamental \emph{degradation mismatch}: standard depth SR assumes bicubic downsampling with optional additive noise, while 3DGS rendering introduces a qualitatively different degradation---splat boundary artifacts, opacity-weighted blending, and coarse quantization.

In this paper, we study depth SR tailored to the 3DGS rendering pipeline.
Our contributions are:

\noindent\textbf{(a) Rendering-induced degradation as a new depth SR setting.}
We show that 3DGS-rendered depth maps exhibit degradation patterns distinct from standard bicubic downsampling.
A 5$\times$5 cross-degradation matrix (\cref{tab:crosstest}) demonstrates a consistent 3--6\,dB PSNR gap between matched and mismatched degradation pairs, establishing that domain-specific training is essential.

\noindent\textbf{(b) Zone-aware evaluation.}
3DGS face models hallucinate geometry in side and back views.
Full-view metrics conflate SR quality with upstream 3DGS artifacts.
We propose restricting evaluation to a near-frontal confidence zone ($|\text{yaw}| \leq 20^\circ$), yielding metrics that accurately attribute SR performance.

\noindent\textbf{(c) 3DGS normals are unusable; DSINE pseudo-GT works.}
Per-splat normal aliasing in 3DGS renders makes surface normals unreliable for normal-aware depth SR losses.
We show that DSINE~\cite{bae2024dsine} monocular normal estimation provides clean pseudo-ground-truth that matches baseline depth SR performance without degradation.

% ===================================================================
\section{Related Work}
\label{sec:related}

\paragraph{Depth super-resolution.}
Early depth SR methods use hand-crafted filters guided by RGB images~\cite{kopf2007joint}.
Learning-based approaches~\cite{hui2016duf,tang2021bridgenet} apply deep networks trained on synthetic LR/HR pairs, typically generated by bicubic downsampling.
Recent methods address degradation robustness: DORNet~\cite{yan2025dornet} learns degradation representations via mixture-of-experts, while SGNet~\cite{yan2024sgnet} uses RGB structure guidance with frequency-domain losses.
SwinIR~\cite{liang2021swinir} and EDSR~\cite{lim2017edsr} are representative single-image SR architectures.
However, none of these methods account for the rendering-specific degradation introduced by 3DGS.

\paragraph{3D Gaussian Splatting.}
3DGS~\cite{kerbl20233dgs} represents scenes as collections of anisotropic 3D Gaussians, enabling real-time rendering via differentiable rasterization.
Extensions include 2DGS~\cite{huang20242dgs} for improved surface reconstruction, Normal-GS~\cite{jiang2024normalgs} for normal supervision, and SuGaR~\cite{guedon2024sugar} for mesh extraction.
FaceLift~\cite{wu2024facelift} applies 3DGS to single-image face reconstruction.
A known limitation across these methods is per-splat normal aliasing~\cite{jiang2024normalgs,huang20242dgs,dai2024dnsplatter}, which we empirically confirm prevents direct use of rendered normals as SR supervision.

\paragraph{Monocular normal estimation.}
DSINE~\cite{bae2024dsine} and Omnidata~\cite{eftekhar2021omnidata} estimate surface normals from a single RGB image using pre-trained models.
We use DSINE pseudo-GT normals as a replacement for noisy 3DGS-rendered normals in our normal-aware loss.

% ===================================================================
\section{Method}
\label{sec:method}

\subsection{Problem formulation}
Given a low-resolution (LR) 8-bit depth map $D_{\text{LR}} \in [0, 255]^{H/s \times W/s}$ rendered by a 3DGS face model, we aim to recover the high-resolution (HR) 16-bit depth map $D_{\text{HR}} \in [0, 65535]^{H \times W}$ at scale factor $s=4$ (256$\to$1024).
After normalization to $[0,1]$, the task becomes:
\begin{equation}
    \hat{D}_{\text{HR}} = f_\theta(D_{\text{LR}}),
\end{equation}
where $f_\theta$ is a learned mapping that jointly performs spatial upsampling and bit-depth recovery.

\subsection{Data pipeline}
\label{sec:data}

We construct our dataset from FFHQ~\cite{karras2019ffhq} images processed through FaceLift~\cite{wu2024facelift}.
Each image yields a 3DGS model from which we render 1024$\times$1024 depth maps (16-bit) at the canonical frontal viewpoint.
LR inputs are generated by resizing HR maps to 256$\times$256 using area interpolation followed by 8-bit quantization.

\textbf{Face masks.}
We generate binary face masks via morphological operations on the cropped face region, restricting both training loss and evaluation metrics to the face area.
This excludes background regions where depth is undefined.

\textbf{Train/val split.}
1{,}158 training and 129 validation samples (90/10 split, seed 42).

\subsection{Network architecture}

We use a UNet~\cite{ronneberger2015unet} with base channels 32 (7.77M parameters).
The input LR depth is first bicubic-upsampled to the target resolution.
The UNet predicts a bounded residual:
\begin{equation}
    \hat{D} = \text{clamp}\left(D_{\text{bic}} + \tanh(U_\theta(D_{\text{bic}})) \cdot 0.5,\ 0,\ 1\right),
\end{equation}
where $D_{\text{bic}}$ is the bicubic-upsampled input.
This residual formulation ensures the network only needs to learn the correction, not the full depth field.

\subsection{Loss function}

The primary loss is a mask-aware L$_1$ loss with a gradient term:
\begin{equation}
    \mathcal{L} = \frac{\sum_{i} M_i \cdot |\hat{D}_i - D_i^{\text{GT}}|}{\sum_i M_i} + 0.5 \cdot \mathcal{L}_{\text{grad}},
\end{equation}
where $M$ is the binary face mask and $\mathcal{L}_{\text{grad}}$ penalizes differences in spatial gradients $(\nabla_x, \nabla_y)$ between prediction and ground truth.

\textbf{Normal-aware variant.}
Optionally, we add a cosine normal loss using DSINE~\cite{bae2024dsine} pseudo-GT normals:
$\mathcal{L}_{\text{normal}} = 1 - \cos(\hat{n}, n_{\text{DSINE}})$,
weighted at 0.2.

\subsection{Training details}

All models are trained with AdamW ($\text{lr}=10^{-4}$, weight decay $10^{-4}$), cosine annealing, and mixed-precision (FP16) on a single RTX~4070 Laptop (8\,GB).
Batch size 2 with gradient accumulation 4 (effective batch 8).
Training converges in $\sim$100 epochs ($\sim$2 hours).

% ===================================================================
\section{Experiments}
\label{sec:experiments}

\subsection{Rendering-induced degradation}
\label{sec:degradation}

To validate that 3DGS rendering produces a distinct degradation, we train five identical UNets on five LR variants generated from the same HR depth maps: (1)~bicubic downsampling, (2)~bicubic + Gaussian noise, (3)~bicubic + 8-bit quantization, (4)~3DGS rendering, and (5)~3DGS rendering + 8-bit quantization.
Each model is evaluated on all five test sets, yielding a 5$\times$5 PSNR matrix (\cref{tab:crosstest}).

\begin{table}[t]
\centering
\caption{Cross-degradation test (PSNR\,dB). Each row is a model trained on the given degradation; each column is the test degradation. Diagonal (bold) = matched. The 3.5\,dB average gap between diagonal and off-diagonal confirms degradation specificity.}
\label{tab:crosstest}
\small
\setlength{\tabcolsep}{3pt}
\begin{tabular}{@{}lccccc@{}}
\toprule
Train $\backslash$ Test & Bic. & +Noise & +Q8 & Render & Rnd+Q8 \\
\midrule
Bicubic      & \textbf{44.6} & 35.8 & 44.5 & 43.2 & 43.1 \\
Bic.+Noise   & 41.2 & \textbf{40.9} & 41.2 & 40.4 & 40.5 \\
Bic.+Q8      & 44.6 & 36.0 & \textbf{44.5} & 43.3 & 43.3 \\
Render       & 40.5 & 32.7 & 40.4 & \textbf{46.2} & 45.9 \\
Render+Q8    & 40.6 & 33.2 & 40.5 & 45.7 & \textbf{45.5} \\
\bottomrule
\end{tabular}
\vspace{-1em}
\end{table}

Key observations:
(1)~Models trained on rendering degradation achieve the highest PSNR (46.2\,dB) when tested on rendering degradation, but drop to 40.5\,dB on bicubic---a 5.7\,dB gap.
(2)~Conversely, bicubic-trained models score 44.6\,dB on bicubic test but only 43.2\,dB on rendering test.
(3)~Two degradation clusters emerge: \{bicubic, bicubic+Q8\} and \{render, render+Q8\}, with noise as an outlier.
This confirms that rendering degradation constitutes a genuinely different SR setting.

\subsection{Baseline comparison}
\label{sec:baselines}

We compare our UNet against six baselines (\cref{tab:baselines}), all trained on the same 3DGS-degraded data except DORNet (NYU-v2 pretrained, zero-shot transfer).

\begin{table}[t]
\centering
\caption{Depth SR results on 129 validation samples (masked metrics, face region only). $^\dagger$SGNet evaluated at 512 resolution due to OOM at 1024.}
\label{tab:baselines}
\small
\setlength{\tabcolsep}{4pt}
\begin{tabular}{@{}llccc@{}}
\toprule
Method & Params & Train & L$_1$ $\downarrow$ & PSNR $\uparrow$ \\
\midrule
Bicubic (no model) & -- & -- & 0.00332 & 41.7 \\
\midrule
SwinIR-tiny~\cite{liang2021swinir} & 0.23M & 3DGS & 0.08561 & 11.9 \\
EDSR-light~\cite{lim2017edsr} & 1.52M & 3DGS & 0.00241 & 46.0 \\
SRResNet~\cite{ledig2017srresnet} & 1.53M & 3DGS & 0.00228 & 46.5 \\
DORNet~\cite{yan2025dornet} & 3.05M & NYU & 0.00337 & 42.1 \\
SGNet$^\dagger$~\cite{yan2024sgnet} & 9.22M & 3DGS & 0.00352 & 43.1 \\
\midrule
UNet (ours) & 7.77M & 3DGS & 0.00232 & 46.2 \\
UNet+Normal (ours) & 7.78M & 3DGS & \textbf{0.00228} & \textbf{46.5} \\
\bottomrule
\end{tabular}
\vspace{-1em}
\end{table}

\textbf{SwinIR-tiny} (0.23M) performs extremely poorly (L$_1$=0.086, PSNR=11.9\,dB), worse than bicubic interpolation.
SwinIR lacks a global skip connection: it must learn the full depth field from scratch rather than predicting a residual.
This confirms the importance of the bicubic residual prior for depth SR.

\textbf{DORNet} (CVPR 2025 oral) achieves only 42.1\,dB---below even simple bicubic interpolation---despite being the current SOTA on NYU-v2.
This 4.1\,dB gap to our method directly demonstrates degradation mismatch: a model trained on bicubic degradation cannot transfer to 3DGS rendering degradation.

\textbf{SGNet} (AAAI 2024, RGB-guided) achieves 43.1\,dB at 512 resolution with RGB guidance.
Despite having access to the RGB image, it underperforms our depth-only UNet by 3.1\,dB, suggesting that the rendering degradation modeling is more important than RGB guidance for this task.

\textbf{SRResNet and EDSR-light}, both equipped with bicubic skip connections and trained on 3DGS data, perform comparably to our UNet (46.0--46.5\,dB), confirming that the residual learning formulation and matched degradation are the primary performance drivers.

\subsection{Normal supervision ablation}
\label{sec:normal}

\begin{table}[h]
\centering
\caption{Normal supervision ablation. 3DGS normals degrade performance due to per-splat aliasing; DSINE pseudo-GT matches the no-normal baseline.}
\label{tab:normal}
\small
\begin{tabular}{@{}lcc@{}}
\toprule
Normal source & Val L$_1$ $\downarrow$ & Note \\
\midrule
None (baseline) & 0.00232 & -- \\
3DGS rendered + bilateral & 0.00370 & Per-splat aliasing \\
DSINE pseudo-GT & \textbf{0.00228} & Clean normals \\
\bottomrule
\end{tabular}
\vspace{-1em}
\end{table}

\cref{tab:normal} shows that using 3DGS-rendered normals as supervision \emph{hurts} performance (L$_1$ increases from 0.00232 to 0.00370), even after bilateral filtering.
This is consistent with the per-splat aliasing documented in Normal-GS~\cite{jiang2024normalgs} and 2DGS~\cite{huang20242dgs}: each Gaussian contributes an independent normal, creating high-frequency artifacts on smooth skin surfaces.

Replacing 3DGS normals with DSINE~\cite{bae2024dsine} pseudo-GT recovers clean supervision, achieving L$_1$=0.00228---matching the best baseline.
This confirms that normal supervision can work for depth SR when the normals are reliable, and that 3DGS-rendered normals are fundamentally unsuitable.

\subsection{Quantization gap}
\label{sec:quant}

\begin{table}[h]
\centering
\caption{Quantization gap: 8-bit \vs 16-bit LR input.}
\label{tab:quant}
\small
\begin{tabular}{@{}lcc@{}}
\toprule
LR input & Val L$_1$ $\downarrow$ & Task \\
\midrule
16-bit (256$\times$256) & 0.00207 & Spatial SR only \\
8-bit (256$\times$256) & 0.00232 & SR + dequantization \\
\midrule
\multicolumn{2}{@{}l}{Gap: 0.00025 (12.1\%)} & \\
\bottomrule
\end{tabular}
\vspace{-1em}
\end{table}

To isolate the dequantization component, we train an identical UNet on 16-bit LR inputs (\cref{tab:quant}).
The 16-bit model achieves L$_1$=0.00207 \vs 0.00232 for 8-bit---a 12.1\% relative gap.
This quantifies the difficulty added by 8-bit quantization and validates our combined SR + dequantization formulation.

% ===================================================================
\section{Conclusion}
\label{sec:conclusion}

We have shown that depth super-resolution for 3DGS face reconstruction requires domain-specific training: the rendering-induced degradation of 3DGS is fundamentally different from the bicubic downsampling assumed by standard depth SR methods.
A simple UNet with bicubic residual learning, trained on 3DGS-rendered data, outperforms the CVPR~2025 SOTA (DORNet) by 4.1\,dB and AAAI~2024 SGNet by 3.1\,dB.
We further demonstrate that zone-aware evaluation is essential to avoid conflating SR quality with upstream 3DGS artifacts, and that DSINE pseudo-GT normals can replace the noisy 3DGS-rendered normals as a supervision signal.

Future work includes feeding SR depth back into 3DGS reconstruction to measure downstream NVS quality improvements, and scaling to $\times$8 SR (128$\to$1024).

% ===================================================================

{\small
\bibliographystyle{ieee_fullname}
\bibliography{references}
}

\end{document}
